\chapter{Introducción}

\section{Motivación}
Las operaciones de Búsqueda y Rescate en entornos naturales o semi-naturales (WiSAR, por sus siglas en inglés: \textit{Wilderness Search and Rescue}) representan misiones críticas de alta complejidad logística en las que el factor tiempo determina directamente la probabilidad de supervivencia del sujeto desaparecido. La literatura médica y operativa en emergencias recalca de forma unánime que las primeras horas de una desaparición (conocidas habitualmente como la \textit{"Golden Hour"} o la ventana dorada de supervivencia) son esenciales, tras este umbral inicial, la probabilidad de encontrar a la víctima con vida decae de manera no lineal y precipitada debido a factores ambientales como la deshidratación, la hipotermia o posibles lesiones físicas sufridas \cite{lyu2023, ewers2025deep, mishra2020drone}.

Tradicionalmente, la localización de personas extraviadas ha dependido de la movilización masiva de recursos terrestres y aéreos. Esta metodología clásica abarca desde patrullas de rescatistas a pie y unidades caninas, hasta helicópteros tripulados de rescate. Aunque efectivos, estos despliegues conllevan notables limitaciones operativas: costes socio-económicos prohibitivos, largos tiempos de organización y un riesgo sustancial para el propio personal de rescate, que con frecuencia debe navegar terrenos accidentados bajo condiciones meteorológicas adversas.

En la última década, la irrupción de los Vehículos Aéreos No Tripulados (UAVs o drones) ha supuesto una revolución tecnológica en la gestión de emergencias civiles, respaldada por proyectos europeos pioneros como ICARUS \cite{icarus2016project} y avances recientes en sensores y visión embarcada para detección de personas en exteriores \cite{alonso2025utilizacion}. Gracias a su flexibilidad operativa, capacidad de vuelo a baja altura y compatibilidad con sensores ópticos y térmicos avanzados, los drones se postulan como el recurso idóneo para agilizar las fases iniciales de una búsqueda \cite{rahman2025}. Sin embargo, la madurez del hardware contrasta con las limitaciones actuales del software de guiado y planificación táctica. La gran mayoría de los despliegues reales de UAVs se limitan a seguir patrones de vuelo estrictamente geométricos y estáticos (como el clásico patrón de cortacésped o \textit{Lawnmower}), o bien dependen del pilotaje manual directo.

Estas aproximaciones convencionales ignoran dos variables cruciales para la optimización táctica:
\begin{enumerate}
    \item \textbf{La asimetría de la probabilidad de presencia (POA):} Una persona perdida no se distribuye de manera uniforme o gaussiana sobre el terreno. Sus movimientos están fuertemente condicionados por la topografía, la presencia de barreras físicas (ríos, pendientes) y vías de comunicación (senderos, carreteras), así como por su propio perfil demográfico y psicológico \cite{koester2008lost}.
    \item \textbf{La limitación energética de los drones:} La autonomía de batería de un UAV multirrotor estándar oscila entre los 20 y 40 minutos. Por lo tanto, un barrido ciego y exhaustivo de grandes extensiones boscosas resulta físicamente inviable, ya que consume gran parte de la energía sobrevolando zonas con nula o bajísima probabilidad de hallazgo.
\end{enumerate}

Para resolver este conflicto, surge la necesidad de formular las operaciones de rescate bajo el prisma de la Búsqueda en Tiempo Mínimo o MTS (\textit{Minimum Time Search}). Este enfoque matemático y algorítmico busca planificar trayectorias para una flota de drones de tal forma que se maximice la probabilidad acumulada de detección a lo largo del tiempo, priorizando la exploración inmediata de los "puntos calientes" de probabilidad antes de proceder al barrido de zonas secundarias \cite{lanillos-bayesian}. La presente investigación se enmarca en este desafío, integrando el modelado de entornos reales mediante SAREnv \cite{drones9090628} con el desarrollo y validación de algoritmos de búsqueda bioinspirados en el framework MTS (\textit{MTS-UncertainEnvironment}) \cite{Broton_Gutierrez_MTS-UncertainEnvironment}. Todo el código fuente, los cuadernos interactivos y los registros experimentales generados en este trabajo se encuentran formalmente integrados en el repositorio del grupo de investigación (rama oficial \texttt{sarenv-mts})\footnote{Repositorio del proyecto: \url{https://github.com/Jompy-GitHub/MTS-UncertainEnvironment/tree/sarenv-mts}. El repositorio se encuentra bajo control de acceso del laboratorio.}.

\subsection{Planteamiento del Problema y Propuesta de Valor de \texttt{sarenv-mts}}

A partir de este contexto operativo, la propuesta desarrollada en este trabajo aborda de forma directa las limitaciones de los métodos tradicionales de búsqueda mediante una formulación rigurosa de ingeniería:

\begin{itemize}
    \item \textbf{¿Qué problema resuelve el framework?:} En una emergencia real de salvamento, los drones multirrotores disponen de una autonomía severamente acotada por la capacidad de sus baterías (entre 20 y 40 minutos de vuelo útil). Los procedimientos estándar de la industria continúan empleando barridos geométricos estáticos ``a ciegas'' (como pasadas paralelas tipo \textit{Lawnmower} o espirales), malgastando el escaso presupuesto energético y temporal en sobrevolar masivamente extensiones de terreno donde la probabilidad de hallar a la persona es prácticamente nula.
    
    \item \textbf{¿Qué aporta exactamente la propuesta \texttt{sarenv-mts}?:} La arquitectura integrada conecta el modelado geoespacial abierto con motores de optimización avanzada para resolver dicha limitación mediante cuatro pilares:
    \begin{enumerate}
        \item \textbf{Escenario probabilístico e inteligente (SAREnv + LPB):} Descarga cartografía vectorial real de OpenStreetMap (caminos, bosques, edificios, ríos) y la transforma en un mapa de creencias o mapa de calor bayesiano $b(v^0)$, ponderando las coberturas del suelo según la psicología y el perfil de la persona perdida (Demencia, Autista o Senderista) a partir de los modelos estadísticos internacionales de Robert Koester (\textit{Lost Person Behavior}) \cite{koester2008lost}.
        \item \textbf{Filtrado de restricciones físicas mediante indexación espacial (R-Tree):} Aplica una máscara geométrica estricta que anula ($P = 0.0$) la probabilidad en masas de agua profundas y cubiertas de edificaciones cerradas. Esto aporta un realismo operativo esencial: impide que la aeronave pierda tiempo y batería sobrevolando tejados de hormigón que su sensor cenital no puede atravesar o láminas acuáticas incompatibles con la víctima en superficie, forzando al planificador a concentrar los recursos en sendas y claros transitables.
        \item \textbf{Pasarela de integración y middleware UTM:} Proporciona la infraestructura de software en Python para proyectar, transformar y discretizar las coordenadas métricas UTM de la cartografía real en matrices numéricas estandarizadas (JSON y NumPy), directamente procesables por los algoritmos de navegación.
        \item \textbf{Planificación de ruta bioinspirada (MTS):} Despliega algoritmos metaheurísticos de inteligencia de enjambre (ABC, BHA y ACO) que navegan dinámicamente sobre el mapa de creencias bayesiano, calculando trayectorias que priorizan los núcleos de alta probabilidad para minimizar el tiempo de intercepción.
    \end{enumerate}
    
    \item \textbf{El resultado cuantitativo demostrado:} Validado mediante un benchmark experimental de 900 simulaciones de Montecarlo sobre las $1.722\text{ hectáreas}$ ($12.75\text{ km}^2$) de la Casa de Campo de Madrid, el sistema demuestra que la búsqueda guiada informada multiplica entre 3x y 6x la tasa de éxito de localización frente a los barridos geométricos a ciegas tradicionales con un único dron y una sola batería, logrando el contacto visual en tiempos medios de entre 5 y 12 minutos (en pleno cumplimiento de la \textit{Golden Hour}).
\end{itemize}

\section{Objetivos}
El objetivo general de este Trabajo de Fin de Máster es diseñar, integrar y validar experimentalmente un sistema inteligente de guiado multiagente (enjambres de drones) orientado a la planificación de rutas eficientes en misiones de búsqueda y rescate sobre entornos reales, maximizando la probabilidad de localización de la víctima dentro de las restricciones de batería de la flota.

Para la consecución de este fin, se definen los siguientes objetivos específicos:
\begin{itemize}
    \item \textbf{Modelado geográfico y probabilístico del entorno de estudio:} Generar una matriz de probabilidad a priori (POA) realista para el parque de la Casa de Campo en Madrid. Esto exige la descarga de datos de código abierto de OpenStreetMap y su procesamiento multicapa (caminos, ríos, edificios) mediante el framework SAREnv, correlacionando el terreno con los perfiles empíricos de comportamiento de personas perdidas (LPB).
    \item \textbf{Desarrollo del puente de comunicación e integración técnica:} Diseñar una arquitectura de software limpia para conectar SAREnv con el motor del planificador MTS. Esto implica estructurar los datos del mapa georreferenciado, normalizarlos y dar soporte a la importación dinámica de mapas reales en formato binario por parte de los planificadores.
    \item \textbf{Implementación y calibración de metaheurísticas bioinspiradas:} Adaptar e integrar algoritmos basados en inteligencia colectiva y optimización metaheurística (en particular, la Optimización por Colonia de Hormigas o ACO, y la Colonia de Abejas Artificiales o ABC) para su ejecución en entornos reales inciertos y discretizados en rejillas.
    \item \textbf{Establecimiento de un marco de validación estocástica (Montecarlo):} Diseñar una batería de simulaciones de Montecarlo sembrando ubicaciones estocásticas de víctimas basadas en la densidad de probabilidad del mapa de calor, con el fin de evaluar la consistencia estadística de cada algoritmo.
    \item \textbf{Análisis comparativo cuantitativo:} Evaluar el rendimiento de las metaheurísticas inteligentes propuestas frente a algoritmos clásicos de cobertura ciega (barrido paralelo) utilizando métricas estándar en la literatura académica SAR, tales como la probabilidad acumulada (APOD), la probabilidad descontada en el tiempo (TDPD) y la tasa empírica de descubrimiento (LPDS).
\end{itemize}

\section{Entorno socio-económico}
El desarrollo de sistemas de planificación de misiones de rescate mediante drones autónomos posee un impacto socio-económico de primer orden. En el ámbito social, la optimización algorítmica incide directamente en la tasa de éxito de localización de personas desaparecidas en escenarios críticos, tales como desorientación de ancianos con patologías cognitivas (Alzheimer o demencia), accidentes de deportistas en montaña o extravío de niños. Reducir el tiempo de localización mitiga el sufrimiento de los familiares y, fundamentalmente, previene la pérdida de vidas humanas.

Desde una perspectiva económica y operativa, la automatización y coordinación de flotas de drones alivia la carga de los servicios de emergencia civiles y fuerzas de seguridad del Estado (como la Guardia Civil, Protección Civil y la Unidad Militar de Emergencias o UME). Las búsquedas tradicionales requieren movilizar decenas de personas durante días con el consecuente gasto en dietas, logística y combustible de aeronaves tripuladas. La sustitución o apoyo temprano mediante drones autónomos permite acotar la zona de búsqueda prioritaria en minutos, optimizando el uso de recursos y reduciendo los riesgos a los que se exponen los propios equipos de rescate al adentrarse en barrancos o terrenos de difícil acceso.

\section{Marco regulador}
La viabilidad práctica de desplegar un enjambre de drones autónomos para operaciones SAR en el espacio aéreo nacional está sujeta a una estricta regulación aeronáutica, coordinada a nivel europeo por la EASA (\textit{European Union Aviation Safety Agency}) y supervisada en España por la AESA (\textit{Agencia Estatal de Seguridad Aérea}). El marco operativo del presente trabajo se sitúa dentro de las directrices del Reglamento de Ejecución (UE) 2021/664 \cite{eu2021uspace}, que establece el entorno del \textit{U-Space}. El \textit{U-Space} define un conjunto de servicios digitales y procedimientos automatizados diseñados para garantizar el acceso seguro de un gran volumen de aeronaves no tripuladas al espacio aéreo, posibilitando misiones complejas en zonas semiurbanas o forestales colindantes con núcleos de población, como es el caso de la Casa de Campo.

Las misiones autónomas multi-dron propuestas se catalogan bajo la categoría ``Específica'' en la normativa vigente, requiriendo un análisis de riesgos según la metodología SORA (\textit{Specific Operations Risk Assessment}). Las trayectorias generadas por los algoritmos deben, por tanto, integrar restricciones operativas tales como zonas de exclusión de vuelo (\textit{Geofencing}) para proteger infraestructuras críticas o evitar colisiones con otros activos del espacio aéreo, un aspecto que los planificadores inteligentes deben asimilar activamente durante su ejecución en la simulación.

\section{Estructura del documento}
Esta memoria de Trabajo de Fin de Máster se estructura de la siguiente forma para guiar al lector en el proceso de desarrollo y validación:
\begin{itemize}
    \item El \textbf{Capítulo 2} expone el Estado del Arte, recorriendo las bases matemáticas de la teoría de búsqueda bayesiana clásica, los modelos de comportamiento humano ante extravíos (LPB), y revisando de manera analítica los algoritmos bioinspirados y de cobertura de caminos vigentes en la literatura científica de 2023 a 2025.
    \item El \textbf{Capítulo 3} detalla el Modelado del Problema, analizando la formalización teórica del espacio de búsqueda y explicando en profundidad el funcionamiento técnico del framework SAREnv para digitalizar y rasterizar la Casa de Campo.
    \item El \textbf{Capítulo 4} aborda el Desarrollo e Integración del software puente, detallando la arquitectura de datos, la estructura del código de carga e inyección en MTS y la adaptación de los scripts experimentales.
    \item El \textbf{Capítulo 5} describe el Proceso de Experimentación, definiendo el diseño del banco de pruebas, las variables de simulación controladas (número de agentes, batería y algoritmos) y la metodología de Montecarlo.
    \item El \textbf{Capítulo 6} presenta el Análisis de Resultados, contrastando el desempeño táctico de las metaheurísticas bioinspiradas frente a las líneas de base convencionales mediante representaciones visuales y estadísticas de las métricas clave.
    \item Por último, el \textbf{Capítulo 7} reúne las Conclusiones del estudio, destacando los logros alcanzados y sugiriendo las líneas de trabajo futuro para extender el sistema.
\end{itemize}
